% E.215: fresh complete mathematical transcription after the workspace outage.
% Not asserted byte-identical to the pre-outage canonical source.
% The original proof was fully read by root and cohomology_audit; its verifier
% was executed by owner and root before the outage. This transcription uses
% the restored complete proofs E.171--174 and E.209--212 as established inputs.
% No new numerical computation or all-depth result is claimed.

\subsection{Primitive support at five and the exact paired matching gate}

\begin{result}{E.215}{Proved}
{A primitive-five ordinary family and its all-weight paired support test}
There is an ordinary completed pair of infinite support with one integral
all-prime class-two section matching the actual ordinary point at every
support prime in every coefficient different from that prime. Its support
$\Sigma$ can satisfy
\[
 \{2,83,173\}\subset\Sigma,
 \qquad (v_2(A),v_2(B))=(-8,-8),
 \qquad (v_{83}(A),v_{83}(B))=(1,0),
 \qquad (v_{173}(A),v_{173}(B))=(0,1),
 \tag{E.215.1}
\]
and every other support prime $q$ can satisfy
\[
 q\equiv5\pmod8,\qquad q\equiv2\pmod9,
 \qquad \operatorname{ord}_{\mathbb F_5^\times}(q)=4,
 \qquad v_5(q^4-1)=1.
 \tag{E.215.2}
\]
Every assigned valuation is an ordinary integer frozen when its prime is
introduced. In fact every $p\in\Sigma$, including the three private
places, is primitive modulo five and has $v_5(p^4-1)=1$.
No coefficient-place or zero-support point condition is asserted.

At coefficient five, any global lift of a support-matched predecessor
through a weight $w\ge3$ with $4\nmid w$ is again support-matched.
Such a global lift exists by E.209 and B.25. In particular the class-two
section has a support-matched class-three lift and a global class-four
lift, without an automatic class-four matching assertion.

There is an exact finite-parameter description at every remaining weight.
Let $w=4k\ge4$, fix an actual support-matched integral class-$(w-2)$
section, and retain its actual local point comparisons. Put
\[
 b=\operatorname{rank}_{\mathbb Z_5}(\Gamma_{w-1}/\Gamma_w),
 \qquad c=\operatorname{rank}_{\mathbb Z_5}(\Gamma_w/\Gamma_{w+1}),
 \qquad s=1+v_5(w),\quad R=\mathbb Z/5^s.
 \tag{E.215.3}
\]
Then one may explicitly choose an integral two-layer cocycle splitting,
a global reference lift, and the actual local frames so that all
class-$w$ support matches are equivalent to
\[
            (L_p-L_2)\lambda=d_p-d_2
            \quad(p\in\Sigma),\qquad \lambda\in R^b,
 \tag{E.215.4}
\]
where $L_p:R^b\to R^c$ and $d_p\in R^c$ are obtained by localization
in those fixed data. Every solution is realized by an integral global
lift. At most $b+2$ support places, including the place two, suffice to
determine whether the entire displayed system has a solution and to
impose all its equations. This includes non-free row submodules over
$R$. It asserts neither consistency nor an effective stopping rule for
finding a controlling set. At weight four there are exactly $25$
possible parameters and at most four controlling support places.

These are assertions about a nonempty selected ordinary class-two family
and its successive actual lifting fibres. They do not construct a
compatible all-depth section or imply finite support.
\end{result}

\begin{proof}
\emph{Reinitializing the arithmetic construction.}
Use the rational point
\[
              t_*=83/256,\qquad1-t_*=173/256,
 \tag{E.215.5}
\]
and prefixes
\[
              a=83U^8,\qquad b=173V^8,\qquad c=256W^8.
 \tag{E.215.6}
\]
Start with $(a,b,c)=(83,173,256)$ and $M=179!$. The coefficients
are pairwise coprime and $83+173=256$, so $(1,1,1)$ is an actual
unit-coordinate solution of the exponent-$M$ diagonal curve over every
completion and the initial sum error is zero. Retain the actual point
section of E.215.5 in the class-two Hall quotients modulo $2^8$ and
modulo nine. The private valuations are exactly E.215.1.

We verify every change needed in E.172--174 rather than presuming that
an arbitrary rational base can replace theirs. At the odd private
primes the dyadic valuation pairs are
\[
 (v_2(83-1),v_2(83^2-1))=(1,3),\qquad
 (v_2(173-1),v_2(173^2-1))=(2,3).
 \tag{E.215.7}
\]
At a fresh prime in E.215.2 they are $(2,3)$. These are within the
bounds used in E.171. Thus the same modulus $2^8$, the reduction through
$2^5$ to $2^4$ killing the real Bockstein, and the one global
$\mathbb Z/8$ twist suffice. At coefficient three the primes
$2,83,173$ and every fresh prime all have
$v_3(p-1)=0$ and $v_3(p^2-1)=1$. The unchanged modulus-nine argument
therefore gives one global $\mathbb Z/3$ twist. These two choices combine
into one global $\mathbb Z/24$ choice, just as in E.171.

The dyadic denominator valuation is now $-8$ rather than $-6$.
It causes no change in the finite-reference comparison: its coefficient
place is omitted in the dyadic off-diagonal assertion, and at every other
coefficient it is the same fixed integer in the rational reference and
all subsequent prefixes. The old infinity-direction local comparison is
therefore retained. The dyadic Hall exponent and conjugation bounds of
E.171--172 are bounds for the fixed finite quotient itself, not bounds
on this one initial valuation.

The actual finite image proof of E.172 applies with private inertia at
83 and173. Those inertia classes give the two Frattini coordinates of
the point cocycle. Their two-primary commutators with complex conjugation
supply twice those coordinates, and their commutators generate the
geometric derived group. Hence the image over every full cyclotomic
field is $\Phi(Q_2)\times Q_3$, exactly as in E.172. The quadratic
classes are now $83,173,-83\cdot173$. A residue one modulo four and
square at both private primes fixes the required square roots. There
is no quadratic condition at five, because their conductors are prime
to five. At infinity keep the inverse-root mark of $-1$ exactly as in
the extended E.170. The actual class-two arithmetic action and its
finite cusp marks remain cyclotomic; no assertion about an arbitrary
higher arithmetic action is used here.

For a complementary ratio $w=c/b,c/a$, or $-b/a$, a private odd
valuation is $1$ or $-1$. The radical degree and maximal-abelian
argument used in E.154 and E.172 therefore apply unchanged. In
particular, over the sufficiently enlarged cyclotomic field $F$ the
cyclic radical has the same degree $h/2$. Its two- and three-primary
intersections are glued by E.170 and E.128 with the same actual cusp
marks. The remaining primary parts have trivial intersection with the
two fixed section fields. Changing the denominator from $64$ to $256$
introduces no new radical obstruction: its fourth root is rational,
while the private odd valuation still excludes an additional fourth
root in an abelian extension. The negative infinity radical is already
covered by E.170.

\emph{The new five-adic residue factorization.}
The primitive residues with the required minimal shell are precisely
\[
                 \mathcal V=\{2,3,8,12,13,17,22,23\}
                    \subset(\mathbb Z/25)^\times.
 \tag{E.215.8}
\]
For a unit $v$ put $f(v)=(v^4-1)/5\pmod5$. This satisfies
$f(uv)=f(u)+f(v)$. In each fixed primitive residue class modulo five,
its five lifts modulo twenty-five have all five possible values of
$f$: replacing $v$ by $v+5z$ changes $f(v)$ by $4v^3z$, a nonzero
linear coefficient modulo five.

Every target $u\equiv1\pmod5$ is a product of two members of
$\mathcal V$. Choose $a\in\mathbb F_5\setminus\{0,f(u)\}$ and choose
$v\equiv2\pmod5$ with $f(v)=a$. Then $u/v\equiv3\pmod5$ and
$f(u/v)=f(u)-a\ne0$. Both factors belong to $\mathcal V$. Inverse
pairs from $\mathcal V$ pad this to any of the lengths six, eight, or
ten. The allowed conditions are open, so multiplying the last factor
by a sufficiently small principal unit makes the product equal to an
arbitrary prescribed five-adic target in $1+5\mathbb Z_5$, not merely
to its reduction modulo twenty-five.

This is exactly the extra factorization needed in the product-CRT step
of E.172. Newton gives all three local coordinate targets near one at
five, as well as at two and the two private odd primes. At three choose
their signs so they are one modulo three and choose the even counts
six, eight, or ten to give the required residues modulo nine. At two
an even product of residues five modulo eight is one modulo eight.
At each private odd prime the square subgroup, of order 41 or 86,
contains nonexceptional inverse pairs; the target is one in its residue
field. At every other treated prime $r\ge7$, a nonzero target is a
product of two residues outside $\{1,-1\}$ by choosing the first
outside $\{1,-1,u,-u\}$, and inverse pairs pad the length. Principal-unit
adjustments give the required exact local products. All these finite
conditions are combined by CRT.

Consequently the complete arithmetic advance of E.172 holds for the
new base, with E.215.2 imposed at every new prime and with
$M$ divisible by $179!$. Its successive complementary roots, congruences
one modulo earlier newly selected primes, Hensel arguments, and
large-prime diagonal-curve Weil bound are unchanged. The additional
condition at five is a condition on the cyclotomic CRT value, so it
preserves the actual two-/three-primary fixed-root gluing. In particular
every newly added prime remains primitive modulo five with the required
nonzero Fermat digit.

The once-only repair E.173 is now required only at coefficients
$\ell\ge7$. Its first arithmetic Frobenius is chosen with the same
primitive/minimal five-residue requirement. The two private square
conditions are independent of it by CRT, and their square subgroups
still contain nonexceptional elements if the active coefficient equals
83 or 173. The complementary radical and retained two-/three-primary
fields have commutator of order prime to the active $\ell$, as in the
proved E.173. Its global first-coordinate reciprocity condition is
unchanged. Thus that theorem applies literally after these substitutions.
Require its integer correction exponents to be divisible by the
prime-to-$\ell$ part of the current exponent; in particular they remain
divisible by all retained powers of five. There is no once-only repair
at coefficient five: $p^2-1$ is a five-adic unit at every support place,
so its integral fixed-abelian class-two local fibre is already a singleton.

\emph{The completed pair and all-prime comparison.}
Enumerate the coefficients at least seven, each once. At its turn compute
the finite maxima $e,u$ of E.171 from the current support and choose
$N_\ell\ge\max(1,e+u)$. Before the repair, enlarge the ordinary exponent
to a multiple of the current exponent, $(n+179)!$, $\ell^{N_\ell}$,
and all stored moduli. Include $\ell$ among the nonexceptional congruences
in this advance. Its new primes do not increase either maximum, so
this choice of $N_\ell$ is not circular. Apply E.173 once, store its
finite section and support, and thereafter avoid $\pm1$ modulo this
coefficient at every new prime.

For each positive integer $d$, eventually all advances and all later
repairs have exponents divisible by $d$: its prime divisors at least
seven have already had their unique repair, while two, three, and five
have no repair exception. The proof of E.174 now gives convergence in
the full multiplicative completion. At each finite place the valuations
are eventually constant ordinary integers. After removing these fixed
prime powers, the unit sequences converge because the later exponents
kill every prescribed finite local unit group. The increasing sum-error
bound, preserved by every repair, gives actual nonzero limits
$a_{\infty,p}+b_{\infty,p}=c_{\infty,p}$. Hence
$A_p=a_{\infty,p}/c_{\infty,p}$ and
$B_p=b_{\infty,p}/c_{\infty,p}$ are the two coordinates of an actual
ordinary point. At the real place their signs are positive and are those
of the point $1/2$. Each ordinary advance adds at least eighteen new
support primes, so the final support is infinite.

For two and three the retained finite reference classes continue to
match the final actual support points: the valuations are fixed, the
local unit groups are killed by the stored powers, and the construction
retains both the cusp Frobenius and the unit part. At a processed
$\ell\ge7$, its stored old local unit classes and finite abelian
reduction are unchanged, and every later support prime is nonexceptional
at $\ell$. Every hypothesis of E.171 is therefore satisfied. At five
the local group $H^1(\mathbb Q_p,\mathbb Z_5(2))$ is zero for every
$p\in\Sigma$. The unique global integral five-primary class-two class
therefore matches all those points without a correction. E.171's single
$\mathbb Z/24$ twist assembles the coefficients into one integral
all-prime class-two section. This proves the existence assertion and
all its stated support conditions.

\emph{Automatic weights and the actual two-layer module.}
Five satisfies the class-number hypothesis of E.209: the class number of
$\mathbb Q(\mu_5)$ is one, as established by the Minkowski calculation
in E.202. Thus the regular-prime inputs used below apply at five.
The private residues modulo twenty-five are $2,8,23$, all in E.215.8.
For every support prime and every positive weight,
\[
 v_5(p^w-1)=
 \begin{cases}0,&4\nmid w,\\1+v_5(w),&4\mid w.
 \end{cases}
 \tag{E.215.9}
\]
The first line is primitivity, and the second is the binomial valuation
formula applied to $p^4$, whose difference from one has valuation one.
The local integral central correction group in weight $w$ is
$\mathbb Z_5/(p^w-1)$ in each Hall coordinate. Thus at nonmultiples
of four it vanishes. E.209 and B.25 provide a global next lift whenever
the predecessor is support-matched, and every such lift then remains
matched at these weights. This proves the automatic-weight assertion.

Fix now the actual class-$(w-2)$ section in the statement. The group
\[
                 M=\Gamma_{w-1}/\Gamma_{w+1}
 \tag{E.215.10}
\]
is abelian, since $2(w-1)\ge w+1$, and is a finite free
$\mathbb Z_5$ module. Its twisted Galois action is canonically determined
by the fixed predecessor: changing a lift of that predecessor by an
element of $\Gamma_{w-1}$ acts trivially on $M$. It is the actual module
in the exact sequence
\[
 0\longrightarrow L_w\longrightarrow M\longrightarrow L_{w-1}
 \longrightarrow0,
 \qquad L_w=\mathbb Z_5(w)^c,\qquad
 L_{w-1}=\mathbb Z_5(w-1)^b.
 \tag{E.215.11}
\]
No diagonalization of its extension class is assumed.

A global lift through both weights exists: first lift through the odd
weight $w-1$, whose support correction groups vanish, and then apply
E.209 to this matched predecessor. Fix one such class-$w$ lift $t$.
The two-layer global lift classes form its torsor under $H^1(G_{\mathbb Q},M)$.
There is no additional geometric stabilizer: a first nonzero integral
graded coordinate of one would be annihilated by every
$\chi_5(g)^i-1$, impossible in a torsion-free lattice. The same argument
applies to every local predecessor and makes its actual lower aligning
frame unique up to the usual final central coboundary.

The cohomology sequence of E.215.11 gives an injection from
$H^1(G_{\mathbb Q},L_w)$ and a map onto
$H^1(G_{\mathbb Q},L_{w-1})$. This last map is surjective for the
following substantive reason. Every odd-weight global correction changes
the class-$(w-1)$ lift but preserves all support matches, because their
local odd-weight groups are zero. E.209 then lifts that corrected
predecessor through weight $w$. Thus every odd class is the image of a
two-layer class. This does not assume that absolute global $H^2$ vanishes.
Using E.210 gives the exact sequence
\[
 0\longrightarrow R^c\longrightarrow H^1(G_{\mathbb Q},M)
 \longrightarrow\mathbb Z_5^b\longrightarrow0.
 \tag{E.215.12}
\]
Choose continuous integral cocycle lifts of a basis of the free quotient.
Their finite linear combinations give a continuous $\mathbb Z_5$-linear
splitting $j$ of E.215.12. Every global lift is consequently represented
by $t+u+j(\lambda)$, with $u\in R^c$ and $\lambda\in\mathbb Z_5^b$.

Locally, $H^0(\mathbb Q_p,L_{w-1})=0$ and
$H^1(\mathbb Q_p,L_{w-1})=0$, so
\[
 H^1(\mathbb Q_p,L_w)\ \xrightarrow{\ \sim\ }\ H^1(\mathbb Q_p,M).
 \tag{E.215.13}
\]
The localization of the global generator
$(\chi_5^w-1)/5^s$ has value $(p^w-1)/5^s$, a five-adic unit,
in the local group of order $5^s$. Hence global central localization
$R^c\to H^1(\mathbb Q_p,M)$ is an isomorphism at every support place.
Use its inverse as the local coordinate identification. Let $d_p$ be
the correction required by the actual point in its retained frame, and
let $L_p$ be localization of $j$ in these coordinates. It factors through
$R^b$. The full local condition is exactly
\[
                       u+L_p\lambda=d_p.
 \tag{E.215.14}
\]
The place two pins $u=d_2-L_2\lambda$, leaving E.215.4. Every vector
in $R^b$ has an integral lift, so every solution is genuinely realized.

\emph{A finite controlling set without a freeness assumption.}
Adjoin the right-hand side to each scalar row of E.215.4 and let $H$ be
the submodule of $R^{b+1}$ generated by all these augmented rows.
Since $R$ is a finite principal ideal local ring, Smith reduction shows
that any submodule of $R^{b+1}$ has at most $b+1$ generators. Equivalently
$\dim_{\mathbb F_5}H/5H\le b+1$, even when $H$ is not free. The images
of the original augmented rows span this quotient, so choose at most
$b+1$ original rows whose images form a basis. Nakayama's lemma shows
that those actual rows generate $H$. Satisfying them therefore satisfies
every augmented row, by linear evaluation against $(\lambda,-1)$.
They occur at at most $b+1$ support places; add the place two used to
pin $u$. This proves the bound $b+2$. At $w=4$, the free Lie ranks are
$b=2,c=3$ and $s=1$, giving $25$ possible vectors and at most four
controlling places. The argument supplies no bound for their sizes or
an effective recognition of the final row span, and it supplies no
consistency assertion. This proves the theorem with precisely the
stated finite-depth and fixed-predecessor scope.
\end{proof}

\paragraph{Pre-outage finite certificate and reconstruction scope.}
The original verifier was \path{code/ghost_audit_primitive_five.py}, with
its saved JSON and text outputs. Owner and root independently executed
it before the outage. The reported output was PASS: 15 local residue
factorizations, 128 sample primes at most 2000, and 25472 weight checks
for weights two through 200. The original JSON file SHA-256 was
\path{2f6d3eb206d9a74726fb22f22906ecf311acc952ac7514a5d73ae6f47c6730e6};
its payload SHA-256 was
\path{8f1ba9dd56950c0b52a13c0060a76a0bacdf6544d0fca8c2a3857f458108402f}.
These are retained execution records, not a claim that this reconstructed
TeX is byte-identical to the original source or that its verifier was
rerun after the outage. The theorem is proved above; the finite certificate
checks finite ingredients and does not execute the Chebotarev recursion.
